\documentclass[xelatex,aspectratio=169]{beamer}

\hfuzz=10pt
\vfuzz=10pt

% Theme
\usetheme{htw}
\setbeamertemplate{navigation symbols}{}
\setbeamertemplate{theorems}[numbered]
\setbeamercovered{transparent}

%\logo{\includegraphics[height=0.5cm]{HTWD_color.png}}

% Packages
\usepackage{polyglossia}
\setmainlanguage{german}
\setotherlanguage{english}

\usepackage[bigfiles]{pdfbase}
\ExplSyntaxOn
\NewDocumentCommand\embedvideo{smm}{
\group_begin:
\leavevmode
\tl_if_exist:cTF{file_\file_mdfive_hash:n{#3}}{
  \tl_set_eq:Nc\video{file_\file_mdfive_hash:n{#3}}
}{
  \IfFileExists{#3}{}{\GenericError{}{File~`#3'~not~found}{}{}}
  \pbs_pdfobj:nnn{}{fstream}{{}{#3}}
  \pbs_pdfobj:nnn{}{dict}{
    /Type/Filespec/F~(#3)/UF~(#3)
    /EF~<</F~\pbs_pdflastobj:>>
  }
  \tl_set:Nx\video{\pbs_pdflastobj:}
  \tl_gset_eq:cN{file_\file_mdfive_hash:n{#3}}\video
}
%
\pbs_pdfobj:nnn{}{dict}{
  /Type/RichMediaInstance/Subtype/Video
  /Asset~\video
  /Params~<</FlashVars (
  source=#3&
  skin=SkinOverAllNoFullNoCaption.swf&
  skinAutoHide=true&
  skinBackgroundColor=0x5F5F5F&
  skinBackgroundAlpha=0.75
  )>>
}
%
\pbs_pdfobj:nnn{}{dict}{
/Type/RichMediaConfiguration/Subtype/Video
/Instances~[\pbs_pdflastobj:]
}
%
\pbs_pdfobj:nnn{}{dict}{
/Type/RichMediaContent
/Assets~<<
/Names~[(#3)~\video]
>>
/Configurations~[\pbs_pdflastobj:]
}
\tl_set:Nx\rmcontent{\pbs_pdflastobj:}
%
\pbs_pdfobj:nnn{}{dict}{
  /Activation~<<
  /Condition/\IfBooleanTF{#1}{PV}{XA}
  /Presentation~<</Style/Embedded>>
  >>
  /Deactivation~<</Condition/PI>>
}
%
\hbox_set:Nn\l_tmpa_box{#2}
\tl_set:Nx\l_box_wd_tl{\dim_use:N\box_wd:N\l_tmpa_box}
\tl_set:Nx\l_box_ht_tl{\dim_use:N\box_ht:N\l_tmpa_box}
\tl_set:Nx\l_box_dp_tl{\dim_use:N\box_dp:N\l_tmpa_box}
\pbs_pdfxform:nnnnn{1}{1}{}{}{\l_tmpa_box}
%
\pbs_pdfannot:nnnn{\l_box_wd_tl}{\l_box_ht_tl}{\l_box_dp_tl}{
  /Subtype/RichMedia
  /BS~<</W~0/S/S>>
  /Contents~(embedded~video~file:#3)
  /NM~(rma:#3)
  /AP~<</N~\pbs_pdflastxform:>>
  /RichMediaSettings~\pbs_pdflastobj:
  /RichMediaContent~\rmcontent
}
\phantom{#2}
\group_end:
}
\ExplSyntaxOff


\usepackage{graphicx}
\usepackage[export]{adjustbox}
\usepackage{animate}
%\usepackage[dvipdfmx]{movie15_dvipdfmx}
\usepackage{media9}
\usepackage{tabularx}
\usepackage{colortbl}
\usepackage{booktabs}
\usepackage{makecell}
\usepackage{ltablex}
\usepackage{array}
\usepackage{multirow}
\usepackage{amsmath}
\usepackage{amsthm}
%\renewcommand{\arraystretch}{1.5}
\newcolumntype{L}[1]{>{\raggedright\let\newline\\\arraybackslash\hspace{0pt}}p{#1}}
\newcolumntype{C}[1]{>{\centering\let\newline\\\arraybackslash\hspace{0pt}}p{#1}}
\newcolumntype{R}[1]{>{\raggedleft\let\newline\\\arraybackslash\hspace{0pt}}p{#1}}
%\renewcommand\thesatz{\arabic{section}.\arabic{theorem}}
\makeatletter
\@addtoreset{theorem}{lecture}
\makeatother

\newtheorem{satz}{Satz}[section]
\newtheorem{lem}{Lemma}[section]
\newtheorem{beh}{Behauptung}[section]
\newtheorem{define}{Definition}[section]
\numberwithin{equation}{section}
\usepackage{ragged2e}
\usepackage{etoolbox}

\usepackage{color}
\usepackage{colortbl}
\definecolor{hellgrau}{rgb}{0.85,0.85,0.85}
\definecolor{hellrot}{rgb}{1,0.7,0.7}

\usepackage{tikz}
\usetikzlibrary{shapes,arrows.meta,calc,arrows,positioning,patterns,tikzmark}
%\usepackage{tikz-uml}
\usepackage{pgfplots}  % for elliptic curves (part 8)
\pgfplotsset{compat=1.18}
\usepackage{pgffor}
\usepackage{pgfmath-xfp}
\tikzset{>=latex}
\tikzset{
  invisible/.style={opacity=0},
  visible on/.style={alt={#1{}{invisible}}},
  alt/.code args={<#1>#2#3}{%
      \alt<#1>{\pgfkeysalso{#2}}{\pgfkeysalso{#3}} % \pgfkeysalso doesn't change the path
    },
}

\usepackage{paralist}

\usepackage{url}
\def\UrlBreaks{\do\/\do-}
\PassOptionsToPackage{hyphens}{url}\usepackage{hyperref}

\usepackage[normalem]{ulem} % gestrichelte Unterstreichung (\dashuline{})
\usepackage{cancel}

\makeatletter
\renewcommand{\itemize}[1][]{%
  \beamer@ifempty{#1}{}{\def\beamer@defaultospec{#1}}%
  \ifnum \@itemdepth >2\relax\@toodeep\else
    \advance\@itemdepth\@ne
    \beamer@computepref\@itemdepth% sets \beameritemnestingprefix
    \usebeamerfont{itemize/enumerate \beameritemnestingprefix body}%
    \usebeamercolor[fg]{itemize/enumerate \beameritemnestingprefix body}%
    \usebeamertemplate{itemize/enumerate \beameritemnestingprefix body begin}%
    \list
    {\usebeamertemplate{itemize \beameritemnestingprefix item}}
    {\def\makelabel##1{%
        {%
            \hss\llap{{%
                  \usebeamerfont*{itemize \beameritemnestingprefix item}%
                  \usebeamercolor[fg]{itemize \beameritemnestingprefix item}##1}}%
          }%
      }%
    }
  \fi%
  \beamer@cramped%
  \justifying% NEW
  %\raggedright% ORIGINAL
  \beamer@firstlineitemizeunskip%
}
\makeatother

\apptocmd{\frame}{}{\justifying}{}

\renewcommand\theadfont{\bfseries\sffamily}
\usepackage{ragged2e}
\usepackage{newpxtext}

\setsansfont{texgyreheros}[
  Scale=MatchLowercase,
  UprightFont=*-regular,
  BoldFont=*-bold,
  ItalicFont=*-italic,
  BoldItalicFont=*-bolditalic,
]

% Title
\usepackage[usetransparent=false]{svg}
% Import references
\usepackage[backend=biber,style=numeric,sorting=none]{biblatex}
\addbibresource{references.bib}

%\AtBeginSection[]{
%  \begin{frame}
%    \vfill
%    \centering
%    \begin{beamercolorbox}[sep=8pt,center,shadow=true,rounded=true]{title}
%      \usebeamerfont{title}\thesection.~\secname\par%
%    \end{beamercolorbox}
%    \vfill
%  \end{frame}
%}

\makeatletter
\newenvironment{noheadline}{
  \setbeamertemplate{headline}{}
  \addtobeamertemplate{frametitle}{\vspace*{-0.9\baselineskip}}{}
}{}
\makeatother


\usepackage{xcolor}
\usepackage{algorithm}
\usepackage[linesnumbered,ruled,lined,commentsnumbered,algo2e,ngerman,ngermankw]{algorithm2e}
\usepackage{algorithmic}
\usepackage{caption}
\usepackage[newfloat]{minted}
\captionsetup[listing]{position=top}
\definecolor{mintedbg}{HTML}{282828}
\setminted{
  breaklines=true,
  bgcolor=mintedbg,
  style=monokai,
  formatcom=\color{white}
}
\usepackage{etoolbox}
\makeatletter
% replace \medskip before and after the box with nothing, i.e., remove it
\patchcmd{\minted@colorbg}{\medskip}{}{}{}
\patchcmd{\endminted@colorbg}{\medskip}{}{}{}
\makeatother

\renewcommand{\theFancyVerbLine}{\textcolor{black}{\arabic{FancyVerbLine}}}

\usepackage{pifont}
\newcommand{\cmark}{\ding{51}}%
\newcommand{\xmark}{\ding{55}}%

\newenvironment{changemargin}[2]{%
  \begin{list}{}{%
      \setlength{\topsep}{0pt}%
      \setlength{\leftmargin}{#1}%
      \setlength{\rightmargin}{#2}%
      \setlength{\listparindent}{\parindent}%
      \setlength{\itemindent}{\parindent}%
      \setlength{\parsep}{\parskip}%
    }%
    \item[]}{\end{list}}


\usepackage{csquotes}

% Title
\title{Numpy}
\author{Prof. Dr. Lukas Iffländer}
\institute{HTW Dresden}
\date{}
\usepackage{svg}

% Begin document
\begin{document}

% Title slide
\begin{frame}
    \titlepage
\end{frame}

\begin{frame}{Was ist Numpy?}
    Numpy ist eine Python-Bibliothek, die für numerische Berechnungen optimiert ist. Sie bietet:
    \begin{itemize}
        \item Effiziente Arrays und Matrizen
        \item Eine Vielzahl von mathematischen Funktionen
        \item Unterstützung für lineare Algebra, Fourier-Transformationen und Zufallszahlen
        \item Integration mit anderen Bibliotheken wie SciPy und Matplotlib
        \item Hohe Leistung durch Vektorisierung und Broadcasting
    \end{itemize}

\end{frame}

\section{Matritzen und Vektoren}

\begin{frame}{Matritzen und Vektoren}
    \renewcommand{\arraystretch}{0.9}
    \begin{tabular}{cc}
        \( \begin{bmatrix}
               1
           \end{bmatrix} \)                      & Skalar (1x1 Matrix)       \\ \\
        \(\begin{bmatrix}
              0.1 & 1.2 & 3.1 & 2.0 & -0.1 & 0 & 3.9
          \end{bmatrix}\) & Zeilenvektor (1x7 Matrix)                        \\ \\
        \(\begin{bmatrix}
              0.1 \\
              1.2 \\
              3.1 \\
              2.0 \\
              3.9
          \end{bmatrix}\)                       & Spaltenvektor (5x1 Matrix) \\ \\
        \( \begin{bmatrix}
               1 & 2 & 3 \\
               4 & 5 & 6 \\
               7 & 8 & 9
           \end{bmatrix} \)                      & Quadratische Matrix       \\ \\
        \( \begin{bmatrix}
               1 & 2 & 3 & 4 \\
               5 & 6 & 7 & 8 \\
           \end{bmatrix} \)
                                                & 2x4 Matrix
    \end{tabular}

\end{frame}

\begin{frame}{Was kann Numpy mit Matrizen anstelle?}
    Die gängigen Operationen mit Vektoren, Matrizen und skalaren Variablen
    sind in numpy eingebaut. Das sind beispielsweise:
    \begin{itemize}
        \item Transposition
        \item Skalarprodukt
        \item Matrix-Multiplikation (Falk‘sches Schema)
        \item Matrix-Invertierung
    \end{itemize}
    Dokumentation: \url{https://numpy.org/doc/2.2/}
\end{frame}

\begin{frame}{Numpy Arrays}
    Numpy Arrays sind die zentralen Datenstrukturen in Numpy. Sie sind:
    \begin{itemize}
        \item Homogen: Alle Elemente haben den gleichen Datentyp
        \item Mehrdimensional: Unterstützung für Vektoren, Matrizen und höhere Dimensionen
        \item Effizient: Optimiert für numerische Berechnungen
    \end{itemize}
    Erstellen eines Numpy Arrays:
    \inputminted{python}{src/numpy_arrays.py}
\end{frame}

\begin{frame}{Numpy Arrays}{Zuweisung}
    \inputminted[lastline=12]{python}{src/numpy_arrays_example.py}
\end{frame}
\begin{frame}{Numpy Arrays}{Zugriff}
    Der Zugriff erfolgt mittel Bereichsindexierung:

    \inputminted[firstline=14, lastline=24]{python}{src/numpy_arrays_example.py}
\end{frame}
\begin{frame}{Numpy Arrays}{Abmessungen}
    Von den Objekten kann mittels shape-Eigenschaft die Abmessung in den jeweiligen Dimensionen ausgelesen werden.
    \inputminted[firstline=26, lastline=29]{python}{src/numpy_arrays_example.py}
\end{frame}

\begin{frame}{Matrixmultiplikation}
    Numpy unterstützt die Matrixmultiplikation mit dem @-Operator oder der \texttt{dot()}-Funktion. Beispiel:
    \inputminted[firstline=3]{python}{src/numpy_arrays_multiply.py}
\end{frame}
\begin{frame}{Matrixmultiplikation}{Visuelle Darstellung}
    Bei der Matrixmultiplikation wird jede Zeile der ersten Matrix mit jeder Spalte der zweiten Matrix Elementweise multipliziert und die Ergebnisse addiert.
    \vspace{0.5cm}
    \begin{columns}
        \begin{column}{0.4\textwidth}
            \begin{center}
                \flushright
                $\begin{bmatrix}
                        \textcolor{white}{0}2 & \textcolor{white}{0}\textcolor{blue}{0}  & \textcolor{white}{0}0 \\
                        \textcolor{white}{0}0 & \textcolor{white}{0}\textcolor{olive}{2} & \textcolor{white}{0}0 \\
                        \textcolor{white}{0}0 & \textcolor{white}{0}\textcolor{red}{0}   & \textcolor{white}{0}2
                    \end{bmatrix}$ \\
                $\begin{bmatrix}
                        1                   & 2                    & 3                  \\
                        \textcolor{blue}{4} & \textcolor{olive}{5} & \textcolor{red}{6} \\
                        7                   & 8                    & 9
                    \end{bmatrix}$
                $\begin{bmatrix}
                        2  & 4           & 6  \\
                        8  & \textbf{10} & 12 \\
                        14 & 16          & 18
                    \end{bmatrix}$
            \end{center}
        \end{column}
    \end{columns}
    \vspace{0.5cm}
    Beispiel für das mittlere Feld: \(\textcolor{blue}{4} \cdot \textcolor{blue}{0} + \textcolor{olive}{5} \cdot \textcolor{olive}{2} + \textcolor{red}{6} \cdot \textcolor{red}{0} = \textbf{10}\)
\end{frame}

\begin{frame}{Multiplikation zwischen Vektoren}
    \begin{block}{Multiplikation Zeilen- mal Spaltenvektor}
        Hier ergibt sich ein Skalar.
        \inputminted[firstline=3, lastline=6]{python}{src/numpy_arrays_multiply_vectors.py}
    \end{block}
    \begin{block}{Multiplikation Spalten- mal Zeilenvektor}
        Hier ergibt sich eine Matrix.
        \inputminted[firstline=3, lastline=6]{python}{src/numpy_arrays_multiply_vectors.py}
    \end{block}
\end{frame}

\begin{frame}{Transposition}
    Die Transposition einer Matrix oder eines Vektors vertauscht Zeilen und Spalten.
    \begin{block}{Vektortransposition}
        \inputminted[firstline=3]{python}{src/numpy_vector_transpose.py}
    \end{block}
    \begin{block}{Matrixtransposition}
        \inputminted[firstline=3]{python}{src/numpy_matrix_transpose.py}
    \end{block}

\end{frame}

\begin{frame}{Matrixinvertierung}
    Die Inverse einer Matrix ist eine Matrix, die mit der ursprünglichen Matrix multipliziert die Einheitsmatrix ergibt.

    \small
    \begin{block}{Matrixinvertierung}
        \inputminted[firstline=3]{python}{src/numpy_matrix_invert.py}
    \end{block}

\end{frame}

\begin{frame}{Weitere NumPy-Funktionen}
    \begin{itemize}
        \item elementweise Operationen auf Arrays (ähnlich der map-Funktion)
              \inputminted[firstline=3]{python}{src/numpy_element_wise.py}

        \item Lösung linearer Gleichungssysteme \(A \times x = b\)
              \inputminted[firstline=3,lastline=6]{python}{src/numpy_solve.py}
    \end{itemize}

\end{frame}

\end{document}
